\documentclass[10pt,landscape,a4paper]{article}
\usepackage[utf8]{inputenc}
\usepackage[landscape]{geometry}
\usepackage{multicol}
\usepackage{xcolor}
\usepackage{listings}
\usepackage{tcolorbox}
\usepackage{enumitem}
\usepackage{fancyhdr}
\usepackage{titlesec}

% Page setup
\geometry{top=0.5cm,left=0.5cm,right=0.5cm,bottom=0.5cm}
\pagestyle{fancy}
\fancyhf{}
\renewcommand{\headrulewidth}{0pt}
\setlength{\parindent}{0pt}
\setlength{\parskip}{0pt}

% Colors
\definecolor{sectioncolor}{RGB}{231,76,60}
\definecolor{codecolor}{RGB}{44,62,80}
\definecolor{examplecolor}{RGB}{230,126,34}
\definecolor{bgcolor}{RGB}{236,240,241}

% Section formatting
\titleformat{\section}{\large\bfseries\color{sectioncolor}}{}{0em}{}[\titlerule]
\titlespacing*{\section}{0pt}{1ex}{0.5ex}

% Code listing style
\lstset{
    basicstyle=\ttfamily\small,
    backgroundcolor=\color{bgcolor},
    frame=single,
    breaklines=true,
    xleftmargin=2pt,
    xrightmargin=2pt,
    aboveskip=2pt,
    belowskip=2pt
}

\begin{document}

\begin{center}
    {\Huge\bfseries\color{sectioncolor} GDB \& pwndbg Cheat Sheet}\\
    \vspace{0.2cm}
    {\large GNU Debugger for binary exploitation and reverse engineering}
\end{center}

\begin{multicols}{3}
\setlength{\premulticols}{1pt}
\setlength{\postmulticols}{1pt}
\setlength{\multicolsep}{1pt}
\setlength{\columnsep}{1.5em}

\section*{Starting GDB}
\begin{lstlisting}
# Start with binary
gdb ./binary
gdb -q ./binary          # Quiet mode

# With arguments
gdb --args ./bin arg1 arg2

# Attach to process
gdb -p <pid>
gdb attach <pid>

# Run with input
gdb ./binary
(gdb) run < input.txt
(gdb) run $(python -c 'print("A"*100)')
\end{lstlisting}

\section*{Basic Commands}
\begin{lstlisting}
# Running
run [args]               # Start program
r                        # Short form
start                    # Run and break at main
continue                 # Continue execution
c                        # Short form
kill                     # Kill process
quit / q                 # Exit GDB

# Stepping
next / n                 # Step over (next line)
step / s                 # Step into (enter function)
nexti / ni               # Step over instruction
stepi / si               # Step into instruction
finish                   # Run until return
\end{lstlisting}

\section*{Breakpoints}
\begin{lstlisting}
# Set breakpoints
break main               # Break at function
b *0x08048400           # Break at address
break *main+10          # Offset from symbol
break file.c:42         # Break at line
break vuln if $eax==0   # Conditional

# Manage breakpoints
info breakpoints        # List all
i b                     # Short form
delete 1                # Delete breakpoint 1
delete                  # Delete all
disable 1               # Disable breakpoint
enable 1                # Enable breakpoint
clear function          # Clear at function
\end{lstlisting}

\section*{Examining Memory}
\begin{lstlisting}
# Examine memory
x/NFU address           # Examine N items
# N = count, F = format, U = unit

# Format (F):
x/x    # Hexadecimal
x/d    # Decimal
x/u    # Unsigned decimal
x/o    # Octal
x/t    # Binary
x/c    # Character
x/s    # String
x/i    # Instruction (disassemble)

# Unit size (U):
b      # Byte (1 byte)
h      # Halfword (2 bytes)
w      # Word (4 bytes)
g      # Giant (8 bytes)

# Examples
x/20x $esp              # 20 hex words at ESP
x/10i $eip              # 10 instructions at EIP
x/s 0x08048000          # String at address
x/20gx $rsp             # 20 qwords (64-bit)
\end{lstlisting}

\section*{Registers (32-bit)}
\begin{lstlisting}
# View registers
info registers          # All registers
i r                     # Short form
info registers eax      # Specific register

# Individual registers
p $eax                  # Print EAX
p/x $eip               # Print EIP in hex
set $eax = 0           # Set EAX to 0

# Common 32-bit registers
$eax                    # Accumulator
$ebx                    # Base
$ecx                    # Counter
$edx                    # Data
$esi                    # Source index
$edi                    # Destination index
$ebp                    # Base pointer
$esp                    # Stack pointer
$eip                    # Instruction pointer
\end{lstlisting}

\section*{Registers (64-bit)}
\begin{lstlisting}
# 64-bit registers
$rax, $rbx, $rcx, $rdx
$rsi, $rdi              # Args 1 & 2
$rbp, $rsp, $rip
$r8 - $r15              # Additional regs

# EFLAGS
$eflags
info registers eflags
\end{lstlisting}

\section*{Stack Operations}
\begin{lstlisting}
# View stack
x/20wx $esp             # 20 words from ESP
x/20gx $rsp             # 20 qwords from RSP
info frame              # Current frame info
backtrace / bt          # Call stack
bt full                 # With local vars

# Stack frame
frame 0                 # Select frame
up / down               # Move up/down frames
info args               # Function arguments
info locals             # Local variables
\end{lstlisting}

\section*{Disassembly}
\begin{lstlisting}
# Disassemble
disassemble main        # Disasm function
disas main              # Short form
disas                   # Current function
disas 0x08048400        # At address
disas main, main+50     # Range

# Show context
set disassembly-flavor intel
disas /r main           # With raw bytes
x/10i $eip              # Next 10 instructions
\end{lstlisting}

\section*{Searching Memory}
\begin{lstlisting}
# Find patterns
find /b 0x08048000, 0x08049000, 0x41
find &system, +9999999, "/bin/sh"
search-pattern "/bin/sh"

# Search in sections
search-pattern "flag{" readable

# grep for strings
grep flag
\end{lstlisting}

\section*{Info Commands}
\begin{lstlisting}
# Information
info functions          # List functions
info variables          # List variables
info breakpoints        # List breakpoints
info registers          # List registers
info frame              # Stack frame info
info proc mappings      # Memory mappings
info files              # Files and sections
info symbol 0x8048000   # Symbol at address

# Symbols
p system                # Address of system
p &system
\end{lstlisting}

\section*{pwndbg Commands}
\begin{lstlisting}
# Context view
context                 # Show context
context reg             # Only registers
context stack           # Only stack
context code            # Only disassembly

# Memory search
search "/bin/sh"
search -t string "flag"
search -t pointer 0x08048000

# ROP gadgets
ropper                  # Find ROP gadgets
ropgadget

# Cyclic pattern
cyclic 200              # Generate pattern
cyclic -l 0x61616162    # Find offset

# Heap
heap                    # Heap info
bins                    # Show bins
vis_heap_chunks         # Visualize heap

# ELF info
checksec                # Security features
vmmap                   # Memory mappings
elf                     # ELF header info
got                     # GOT table
plt                     # PLT table
\end{lstlisting}

\section*{Common pwndbg Shortcuts}
\begin{lstlisting}
# Navigation
n / ni                  # Next / Next instruction
s / si                  # Step / Step instruction
c                       # Continue
fin                     # Finish function

# Views
ctx                     # Context
regs                    # Registers
stack                   # Stack view
code                    # Code/disassembly

# Breakpoints
b *main                 # Break at main
b *0x08048400          # Break at address
\end{lstlisting}

\section*{Setting Values}
\begin{lstlisting}
# Set register
set $eax = 0x41414141
set $rip = 0x400000

# Set memory
set {int}0x08048400 = 0x90909090
set {long}$rsp = 0xdeadbeef

# Set string
set {char[4]}0x8048400 = "AAAA"
\end{lstlisting}

\section*{Conditional Breakpoints}
\begin{lstlisting}
# Break with condition
break *main if $eax == 0
break vuln if *(int*)($esp+4) == 0x41

# Commands on break
break *vuln
commands
  x/10wx $esp
  continue
end
\end{lstlisting}

\section*{Working with Core Dumps}
\begin{lstlisting}
# Generate core dump
ulimit -c unlimited     # Enable core dumps
gdb ./binary core       # Load core dump

# After crash
bt                      # Backtrace
info registers          # See crash state
x/10wx $esp            # See stack
\end{lstlisting}

\section*{Scripting}
\begin{lstlisting}
# GDB script (commands.gdb)
b *main
commands
  x/10i $eip
  continue
end
run

# Run script
gdb -x commands.gdb ./binary

# Python scripting
python print("Hello from Python")
python gdb.execute("info registers")
\end{lstlisting}

\section*{Useful Tricks}
\begin{lstlisting}
# Follow fork
set follow-fork-mode child
set follow-fork-mode parent

# Disable ASLR
set disable-randomization on

# Show syscalls
catch syscall
catch syscall write

# Layout (TUI mode)
layout asm              # Assembly view
layout regs             # Registers view
layout split            # Split view
Ctrl-x a                # Toggle TUI

# Redirect output
set logging on
set logging file gdb.txt
\end{lstlisting}

\section*{Pattern Matching}
\begin{lstlisting}
# Generate pattern (pwndbg)
cyclic 200
pwndbg> cyclic 200

# Find offset
cyclic -l 0x6161616b   # From crash value
cyclic -l kaaa         # From string

# Or use pattern_create (GEF/peda)
pattern create 200
pattern offset 0x6161616b
\end{lstlisting}

\section*{Remote Debugging}
\begin{lstlisting}
# On remote machine
gdbserver :1234 ./binary

# On local machine
gdb ./binary
target remote host:1234
\end{lstlisting}

\section*{Essential .gdbinit}
\begin{lstlisting}
# ~/.gdbinit
set disassembly-flavor intel
set pagination off
set follow-fork-mode child

# Load pwndbg
source ~/pwndbg/gdbinit.py

# Custom aliases
define hook-stop
  info registers
  x/10i $eip
end
\end{lstlisting}

\section*{Quick Reference}
\begin{itemize}[leftmargin=*,noitemsep,topsep=0pt]
    \item \texttt{r} - run program
    \item \texttt{c} - continue execution
    \item \texttt{n} - next line (step over)
    \item \texttt{s} - step into function
    \item \texttt{ni} - next instruction
    \item \texttt{si} - step instruction
    \item \texttt{b *addr} - breakpoint at address
    \item \texttt{x/10wx \$esp} - examine stack
    \item \texttt{i r} - info registers
    \item \texttt{disas func} - disassemble function
    \item \texttt{bt} - backtrace (call stack)
    \item \texttt{vmmap} - memory mappings (pwndbg)
    \item \texttt{checksec} - security features (pwndbg)
    \item \texttt{cyclic 200} - pattern (pwndbg)
    \item \texttt{telescope \$esp} - examine memory (pwndbg)
\end{itemize}

\vfill
\columnbreak

\end{multicols}

\end{document}
